\section{Anhang B: Reproduzierbarkeit und Projektlogik}

\subsection{Projektstruktur}
Das Projekt ist in zwei Hauptkomponenten unterteilt: das Node.js-Backend und das React-Native-Frontend. Die folgende Struktur zeigt die wesentlichen Verzeichnisse.
\begin{lstlisting}[language=bash]
.
|-- Backend/school-access-control-backend/
|   |-- middleware/ (Authentifizierung & Sicherheit)
|   |-- routes/ (API-Endpunkte)
|   |-- utils/ (Bildverarbeitung & KI-Logik)
|   |-- weights/ (TensorFlow.js Modellgewichte)
|   `-- server.js (Haupteinstiegspunkt)
|-- Frontend/SchoolAccessControl/ (Mobile App)
|-- BLL-LaTeX/ (Dokumentation)
`-- database.db (SQLite-Datenbank)
\end{lstlisting}

\subsection{Reproduktionsanleitung}
Um das System in einer Testumgebung zu reproduzieren, sind die folgenden Schritte notwendig.

\subsubsection{Hardware-Voraussetzungen}
\begin{itemize}
  \item Raspberry Pi 4 (mind. 4GB RAM empfohlen für KI-Inferenz)
  \item USB-RFID-Lesegerät (125\,kHz, EM4100-kompatibel)
  \item Kamera-Modul oder USB-Webcam zur Gesichtserfassung
\end{itemize}

\subsubsection{Software-Voraussetzungen und Installation}
\begin{enumerate}
  \item Installation von Node.js 22.x und npm.
  \item Repository klonen: \texttt{git clone <repository-url>}.
  \item Backend-Abhängigkeiten installieren:
    \begin{lstlisting}[language=bash]
    cd Backend/school-access-control-backend
    npm install
    npm rebuild @tensorflow/tfjs-node --build-from-source
    \end{lstlisting}
  \item Umgebungsvariablen konfigurieren: Datei \texttt{.env} mit \texttt{JWT\_SECRET} und \texttt{PORT=3000} anlegen.
  \item Server starten: \texttt{node server.js}.
\end{enumerate}

\subsubsection{Durchführung von Benchmarks}
Die im \autoref{sec:ergebnisse} (\nameref{sec:ergebnisse}) genannten Latenzwerte können durch die bereitgestellten Benchmark-Skripte verifiziert werden.
\begin{lstlisting}[language=bash]
cd Backend/school-access-control-backend/benchmarks
node benchmark_face.js
\end{lstlisting}
Das Skript gibt die Inferenzzeit für 20 Testbilder aus und speichert die Ergebnisse in \texttt{results/rpi\_face.csv}.

\subsection{Fehlerbehandlung und Plausibilitätschecks}
\begin{itemize}
  \item \textbf{Modell-Ladefehler:} Falls die Gewichte in \texttt{/weights} fehlen, startet der Gesichtserkennungs-Service nicht. Prüfen Sie das \texttt{server.log} auf entsprechende Fehlermeldungen.
  \item \textbf{Speichermangel:} Bei der Verarbeitung vieler HEIC-Bilder auf einem Raspberry Pi kann es zu \enquote{Out of Memory} Fehlern kommen. Das System ist so konfiguriert, dass es nach jeder Bildkonvertierung \texttt{global.gc()} aufruft (erfordert Start mit \texttt{--expose-gc}).
  \item \textbf{Erwartete Ausgabe:} Ein erfolgreicher RFID-Scan liefert ein JSON-Objekt mit \texttt{"valid": true} und der \texttt{photoUrl} des Schülers zurück.
\end{itemize}
